% Раздел статьи 2: общая структура локальной ошибки. Черновик 16.10.2026.
\subsection{Деревья и элементарные дифференциалы}\label{sec:trees}

Напомним, что в~\cite{PrikBarriers} было получено явное представление младших членов локальной ошибки (предложение~1 там), а остальные члены объединялись в остаток $\rho_n=O(h^{\min(3+\alpha,\,2+3\alpha)})$. Опишем полное разложение.

Основную роль в разложении играют полные производные правой части вдоль решения в точке $t_n$:
\begin{equation}\label{eq:atoms}
D_{r,k}=\frac{\dd^k}{\dd t^k}\,\partial_y^rf\bigl(t,y(t)\bigr)\Big|_{t=t_n},\qquad r,k\ge0,
\end{equation}
так что $D_{0,k}=g_k$, $D_{1,0}=J$, $D_{1,1}=J'$; при $r\ge1$ величина $D_{r,k}$ является симметрическим $r$-линейным отображением. Будем называть их \emph{атомами}. В отличие от классической теории, где элементарные дифференциалы строятся из частных производных $f$, здесь удобнее полные производные: значения $y'(t_n),y''(t_n),\dots$ во внутренней точке определяются всей предысторией, а не значением $f$ в этой точке, и атомы оказываются независимыми параметрами задачи (лемма~\ref{lem:free}).

Индексирующими объектами служат обычные корневые деревья~--- те же, что и в теории Бутчера~\cite{Butcher2016,HNW1993}, но с иной ролью вершин. Вершины дерева $T$ разделим на два типа. \emph{Листья} будем называть \emph{временны\'ми} вершинами: каждый лист отвечает одному дифференцированию по времени и приносит множитель $h$. Все остальные вершины, включая корень, будем называть \emph{стадийными}: каждая стадийная вершина отвечает одному интегрированию с ядром $(t-s)^{\alpha-1}/\G(\alpha)$, т.\,е. одному множителю $h^{\alpha}$ в~\eqref{eq:method}. Дерево из одной вершины состоит из одного корня, который считается стадийным. Обозначим через $N(T)$ число стадийных вершин, через $K(T)$~--- число листьев и положим
\begin{equation}\label{eq:order}
|T|=N(T)\,\alpha+K(T).
\end{equation}
Для стадийной вершины $v$ через $k_v$ обозначим число её детей-листьев, через $r_v$~--- число детей-стадий. Стадийную вершину с $r_v=0$ назовём \emph{висячей}; поскольку она не является листом, у неё $k_v\ge1$ (кроме случая дерева из одной вершины).

\begin{definition}\label{def:F}
Элементарный дифференциал дерева определяется рекурсией по корню $\rho$:
\[
F(T)=D_{r_\rho,k_\rho}\bigl(F(T_1),\dots,F(T_{r_\rho})\bigr),
\]
где $T_1,\dots,T_{r_\rho}$~--- поддеревья с корнями в детях-стадиях корня $\rho$; для висячей вершины $F=g_{k_\rho}$, для дерева из одной вершины $F=g_0$. Коэффициент симметрии $\sigma(T)$~--- обычный коэффициент Бутчера: $\sigma(\bullet)=1$, $\sigma(T)=\prod_i\mu_i!\,\sigma(T_i)^{\mu_i}$, где $\mu_i$~--- кратности попарно различных поддеревьев корня (одинаковые листья учитываются, что даёт множитель $k_\rho!$).
\end{definition}

Например, дереву, у которого корень имеет одного ребёнка-стадию, а тот~--- два листа, отвечают $F=Jg_2$, $|T|=2+2\alpha$, $\sigma=2$; это член $C_1$ работы~\cite{PrikBarriers}. Дереву с корнем, имеющим один лист и одного ребёнка-стадию с одним листом, отвечают $F=J'g_1$ и тот же показатель $2+2\alpha$; это член $C_2$. Цепочка из $m+1$ стадийных вершин, последняя из которых имеет один лист, даёт $F=J^mg_1$ и $|T|=1+(m+1)\alpha$~--- цепочку $B_m$.

\paragraph{Веса.} Пусть задана таблица $(c,A,b)$ явного метода. Каждой некорневой стадийной вершине $v$ сопоставим вектор $\varphi(v)\in\mathbb R^s$ рекурсией снизу вверх:
\begin{equation}\label{eq:phi}
\varphi(v)=\begin{cases}
d_{k_v}=Ac^{k_v}-\gamma_{k_v}\,c^{\,k_v+\alpha}, & r_v=0,\\[1mm]
A\Bigl(c^{k_v}\cdot\varphi(u_1)\cdots\varphi(u_{r_v})\Bigr), & r_v\ge1,
\end{cases}
\qquad \gamma_k=\frac{k!}{\G(k+\alpha+1)},
\end{equation}
где $u_1,\dots,u_{r_v}$~--- дети-стадии вершины $v$, а точка обозначает покомпонентное произведение. Вес дерева определяется корнем:
\begin{equation}\label{eq:Psi}
\Psi(T)=\begin{cases}
D_{k_\rho}=b^{\top}c^{k_\rho}-\gamma_{k_\rho}, & r_\rho=0,\\[1mm]
b^{\top}\Bigl(c^{k_\rho}\cdot\varphi(u_1)\cdots\varphi(u_{r_\rho})\Bigr), & r_\rho\ge1.
\end{cases}
\end{equation}
Квадратурные дефекты $d_k$, $D_k$ здесь те же, что и в~\cite{PrikBarriers}, с точностью до множителя: $\gamma_k$~--- значение интеграла~\eqref{eq:beta} при $c=1$.

Отметим характерную особенность весов~\eqref{eq:phi}, \eqref{eq:Psi}: точное решение входит в них \emph{только через висячие вершины}, т.\,е. через дефекты $d_k$ и $D_k$, а каждая внутренняя вершина несёт чисто численный оператор $A\,\mathrm{diag}(c^{k})$. Причина в том, что разложение правой части в каждой стадии ведётся вокруг точного решения; сравнивать с точным значением приходится лишь квадратуру гладкой функции $g$ на текущем шаге, а не результат применения $f$ к неточным стадиям. Раскрывая произведения дефектов, получаем представление веса через бета-интегралы:
\begin{equation}\label{eq:Psi-beta}
\Psi(T)=\sum_{S\subseteq\mathcal L(T)}(-1)^{|S|}\Bigl(\prod_{v\in S}\gamma_{k_v}\Bigr)\,\Phi_S(T),
\end{equation}
где $\mathcal L(T)$~--- множество висячих вершин, а $\Phi_S(T)$~--- обычный численный вес дерева, в котором для $v\in S$ множитель $Ac^{k_v}$ заменён на $c^{k_v+\alpha}$ (а для висячего корня $b^{\top}c^{k}$~--- на единицу). При $\alpha=1$ имеем $\gamma_kc^{k+1}=c^{k+1}/(k+1)$, и условия $\Psi(T)=0$ треугольным преобразованием сводятся к классическим условиям Бутчера $\Phi(T)=1/\gamma(T)$. При $\alpha<1$ дробные степени $c^{k+\alpha}$, попадающие под действие матрицы $A$, через другие условия не выражаются, и такой редукции нет: именно поэтому каждое дерево даёт отдельное условие.

\paragraph{Связь с деревьями Вольтерры.} Правило~\eqref{eq:order} совпадает с правилом $\mathrm{fin}(t)+(1+a)\,\mathrm{int}(t)$ для деревьев Вольтерры Любича~\cite{Lubich1983}, где $a=\alpha-1$~--- показатель ядра $(x-s)^{a}$, а корень также считается внутренней вершиной. Для ядра общего вида $(x-s)^aK(x,s,y)$ концевые вершины деревьев Любича могут нести дополнительный индекс, отвечающий производной ядра по первому аргументу; для уравнения~\eqref{eq:fde} функция $K=f(s,y)/\G(\alpha)$ от $x$ не зависит, такие вершины не возникают, и индексирующее множество сводится к обычным корневым деревьям. Веса Любича также определены рекурсивно, с дефектами квадратуры в концевых вершинах, и теорема~6 работы~\cite{Lubich1983} утверждает \emph{достаточность} условий $\varphi(t)=0$ для заданного порядка локальной ошибки. Алгоритм автоматической генерации условий для уравнений Вольтерры со слабосингулярным ядром предложен Гаррапой~\cite{Garrappa2010}. Новыми в настоящем разделе являются: представление локальной ошибки как тождества с явными элементарными дифференциалами~\eqref{eq:atoms} и оценкой остатка (теорема~\ref{thm:bseries}), лемма о свободе атомов и вытекающая из неё \emph{необходимость} условий (теорема~\ref{thm:necessity})~--- именно она требуется для доказательства барьеров,~--- а также подсчёт числа условий, обрезание по стадийной высоте и описание скалярного случая.

\subsection{Разложение локальной ошибки}

Обозначим через $\mathcal E=\{N\alpha+K\}$ множество показателей~\eqref{eq:order} всех деревьев; оно дискретно, и для любого $P>0$ определено $P^{+}=\min\{x\in\mathcal E:\ x>P\}$.

\begin{theorem}[разложение локальной ошибки]\label{thm:bseries}
Пусть $P>0$, $M=\lfloor P\rfloor+1$, $f\in C^{M}$ в окрестности дуги $\{(t,y(t)):\,t\in[t_n,t_n+\hat ch]\}$, $y\in C^{M}[t_n,t_n+\hat ch]$. Тогда для метода~\eqref{eq:method}, \eqref{eq:rowsum}
\begin{equation}\label{eq:bseries}
\delta_n=\sum_{T:\ |T|\le P}\frac{h^{|T|}}{\sigma(T)}\,\Psi(T)\,F(T)+\rho_n,\qquad |\rho_n|\le K_P\,h^{P^{+}},
\end{equation}
а ошибки стадий $e_i=Y_i-y(t_n+c_ih)$ представляются в виде
\begin{equation}\label{eq:stage-series}
e_i=\sum_{T\ne\bullet:\ |T|\le P}\frac{h^{|T|}}{\sigma(T)}\,\varphi_i(T)\,F(T)+O\bigl(h^{P^{+}}\bigr),
\end{equation}
где $\varphi(T)$~--- вектор~\eqref{eq:phi}, вычисленный для корня дерева $T$. Постоянная $K_P$ зависит от норм производных $f$ порядка не выше $M$ в окрестности дуги, норм $y^{(j)}$, $j\le M$, от $\max|a_{ij}|$, $\max|b_i|$ и от $\hat c$, но не от $h$; при $y\in C^{M}[t_*,T+\hat ch_0]$, $t_*>0$, она равномерна по шагам с $t_n\ge t_*$, $h\le h_0$.
\end{theorem}

Предложение~1 работы~\cite{PrikBarriers} есть частный случай теоремы~\ref{thm:bseries} при $P<\min(3+\alpha,2+3\alpha)$: деревья с $|T|$ ниже этого значения~--- это в точности $D_0,D_1,D_2$, цепочки $B_m$ и два дерева $C_1,C_2$.

Доказательству предпошлём комбинаторное наблюдение. Каждое дерево с корнем, имеющим $k$ листьев и детей-стадий $T_1,\dots,T_r$, будем записывать как $[\,k;\,T_1,\dots,T_r\,]$.

\begin{lemma}[подстановка рядов]\label{lem:subst}
Пусть $\varepsilon=\sum_{T\in\mathcal S}\frac{h^{|T|}}{\sigma(T)}x(T)F(T)$~--- конечная сумма по множеству $\mathcal S$ деревьев с числовыми коэффициентами $x(T)$. Тогда для $r\ge1$, $k\ge0$
\[
\frac{1}{k!\,r!}\,\tau^k\,D_{r,k}\bigl(\varepsilon,\dots,\varepsilon\bigr)
=\sum_{T_1,\dots,T_r\in\mathcal S}{}^{'}\ \frac{h^{|U|-\alpha-k}\,\tau^k}{\sigma(U)}\,x(T_1)\cdots x(T_r)\,F(U),\qquad U=[\,k;\,T_1,\dots,T_r\,],
\]
где штрих означает суммирование по неупорядоченным наборам $\{T_1,\dots,T_r\}$ (мультимножествам).
\end{lemma}
\begin{proof}
Раскрывая полилинейное выражение, получаем сумму по упорядоченным наборам $(T_1,\dots,T_r)$. Мультимножество с кратностями $\mu_i$ встречается $r!/\prod\mu_i!$ раз, и в силу симметрии $D_{r,k}$ все эти слагаемые равны $D_{r,k}(F(T_1),\dots,F(T_r))=F(U)$. Показатель $h$ равен $\sum|T_l|=|U|-\alpha-k$, а коэффициент~--- $\frac{1}{k!\prod\mu_i!\prod\sigma(T_l)}=\frac{1}{\sigma(U)}$ по определению~\ref{def:F}.
\end{proof}

\begin{proof}[Доказательство теоремы~\ref{thm:bseries}]
\emph{Шаг 1: точное соотношение для ошибок стадий.} Вычитая из~\eqref{eq:method} равенство~\eqref{eq:split} при $t=t_n+c_ih$, получаем
\begin{equation}\label{eq:err-exact}
e_i=h^{\alpha}\Bigl[\sum_ja_{ij}\,g(t_n+c_jh)-\frac{1}{\G(\alpha)}\int_0^{c_i}(c_i-\theta)^{\alpha-1}g(t_n+\theta h)\dd\theta\Bigr]
+h^{\alpha}\sum_{j<i}a_{ij}\bigl[f(t_j,y(t_j)+e_j)-g(t_j)\bigr],
\end{equation}
$t_j=t_n+c_jh$, и то же для $\delta_n$ с заменой $a_{ij}$ на $b_j$ и $c_i$ на $1$.

\emph{Шаг 2: квадратурная часть.} По формуле Тейлора $g(t_n+\theta h)=\sum_{k<M}g_k\frac{(\theta h)^k}{k!}+O(h^M)$ равномерно по $\theta\in[0,\hat c]$; с учётом~\eqref{eq:beta} первая скобка в~\eqref{eq:err-exact} равна
\[
\sum_{k<M}\frac{h^k}{k!}\,(d_k)_i\,g_k+O(h^{M}),
\]
причём $d_0=0$ по~\eqref{eq:rowsum}. Слагаемое с номером $k$ умножается на $h^{\alpha}$ и даёт член $h^{k+\alpha}\frac{(d_k)_i}{k!}g_k$ ряда~\eqref{eq:stage-series}, отвечающий висячему дереву $[\,k;\,]$ с $\sigma=k!$. Остаток имеет порядок $h^{M+\alpha}$; поскольку $M+\alpha=|[\,M;\,]|\in\mathcal E$ и $M+\alpha>P$, имеем $M+\alpha\ge P^{+}$.

\emph{Шаг 3: индукция по номеру стадии.} Метод явный, поэтому $e_1=0$ (первая строка $A$ нулевая, $c_1=0$), а $e_i$ выражается через $e_1,\dots,e_{i-1}$. Предположим, что~\eqref{eq:stage-series} доказано для всех $j<i$. В частности, $e_j=O(h^{1+\alpha})$, так как наименьший показатель некорневого дерева равен $1+\alpha$. Положим $R=\lfloor(P-\alpha)/(1+\alpha)\rfloor$. По формуле Тейлора по второму аргументу
\[
f(t_j,y(t_j)+e_j)-g(t_j)=\sum_{r=1}^{R}\frac{1}{r!}\,\partial_y^rf\bigl(t_j,y(t_j)\bigr)\bigl(e_j,\dots,e_j\bigr)+O\bigl(|e_j|^{R+1}\bigr),
\]
и остаток имеет порядок $h^{(R+1)(1+\alpha)}$. После умножения на $h^{\alpha}$ получаем показатель $\alpha+(R+1)(1+\alpha)$~--- это показатель дерева, корень которого имеет $R+1$ висячих детей с одним листом, и он больше $P$, а значит, не меньше $P^{+}$. Здесь достаточно гладкости $f\in C^{M}$: поскольку $R$~--- целое число, меньшее $P$, имеем $R+1\le M$. Далее, по формуле Тейлора по времени,
\[
\partial_y^rf\bigl(t_n+c_jh,\,y(t_n+c_jh)\bigr)=\sum_{k<M-r}\frac{(c_jh)^k}{k!}\,D_{r,k}+O\bigl(h^{M-r}\bigr).
\]
Остаток, умноженный на $h^{\alpha}|e_j|^r$, имеет порядок $h^{\alpha+M-r+r(1+\alpha)}\le h^{M+\alpha}$, т.\,е. снова $O(h^{P^{+}})$. Подставим теперь в оставшиеся главные члены разложения~\eqref{eq:stage-series} для $e_j$. Произведения остатков $O(h^{P^{+}})$ на ограниченные множители порядка $O(h^{(r-1)(1+\alpha)+\alpha})$ дают $O(h^{P^{+}})$. Произведения главных частей по лемме~\ref{lem:subst} дают суммы по деревьям $U=[\,k;\,T_1,\dots,T_r\,]$ вида
\[
h^{\alpha}\sum_{j<i}a_{ij}\,\frac{h^{|U|-\alpha}}{\sigma(U)}\,c_j^{k}\varphi_j(T_1)\cdots\varphi_j(T_r)\,F(U)=\frac{h^{|U|}}{\sigma(U)}\,\varphi_i(U)\,F(U)
\]
в силу~\eqref{eq:phi}. Каждое такое $U$ с $|U|\le P$ входит в~\eqref{eq:stage-series}, а слагаемые с $|U|>P$ имеют, как показатели деревьев, порядок не ниже $h^{P^{+}}$ и переносятся в остаток. Обратно, каждое дерево с нелистовым корнем и $|U|\le P$ возникает ровно один раз: его поддеревья $T_l$ имеют $|T_l|<|U|-\alpha\le P$ и учтены в разложениях $e_j$, а пара $(r,k)$ удовлетворяет $r\le R$ и $k+r<M$. Действительно, $|U|\ge\alpha+k+r(1+\alpha)$, откуда $r(1+\alpha)\le P-\alpha$ и $k+r\le|U|-\alpha<P<M$. Вместе с шагом 2 это доказывает~\eqref{eq:stage-series} для стадии $i$.

\emph{Шаг 4: локальная ошибка.} Соотношение для $\delta_n$ отличается от~\eqref{eq:err-exact} лишь заменой $a_{ij}\to b_j$, $c_i\to1$ и наличием члена $k=0$ в квадратурной части ($D_0$ не обязан обращаться в нуль). Повторяя шаги 2 и 3, получаем~\eqref{eq:bseries} с весами~\eqref{eq:Psi}; дерево из одной вершины даёт член $h^{\alpha}D_0g_0$.

Все постоянные в оценках зависят лишь от перечисленных в формулировке величин; при $t_n\ge t_*$ они ограничены равномерно, поскольку дуги решения лежат в компакте, на котором $f$ и $y$ имеют ограниченные производные.
\end{proof}

\begin{corollary}[обрезание по стадийной высоте]\label{cor:height}
Пусть $\mathrm{ht}(T)$~--- число стадийных вершин на самом длинном пути от корня к листу. Если $\mathrm{ht}(T)>s$, то $\Psi(T)=0$ для любого явного $s$-стадийного метода.
\end{corollary}
\begin{proof}
Обозначим через $\Pi_m\subset\mathbb R^s$ подпространство векторов, у которых первые $m$ компонент равны нулю. Матрица $A$ строго нижнетреугольная, поэтому $A\Pi_m\subseteq\Pi_{m+1}$ и $A\mathbb R^s\subseteq\Pi_1$; умножение на $c^k$ и покомпонентное произведение с любым вектором сохраняют $\Pi_m$. Индукцией по высоте поддерева получаем $\varphi(v)\in\Pi_{h_v}$, где $h_v$~--- высота поддерева с корнем $v$: для висячей вершины $d_k=Ac^k-\gamma_kc^{k+\alpha}\in\Pi_1$, так как $c_1=0$ и первая строка $A$ нулевая; для внутренней вершины $\varphi(v)=A(\dots)$, где в произведении участвует вектор из $\Pi_{h_v-1}$. Если $\mathrm{ht}(T)>s$, то у корня есть ребёнок $u$ с $h_u\ge s$, т.\,е. $\varphi(u)\in\Pi_s=\{0\}$, и $\Psi(T)=0$.
\end{proof}

При $T=$ цепочка это в точности лемма о полной цепочке из~\cite{PrikBarriers}: $B_m=0$ при $m\ge s$.

\subsection{Необходимость условий}

\begin{lemma}[о свободе атомов]\label{lem:free}
Для любых $t_n>0$, $d\ge1$, натуральных $R,K$, векторов $\hat g_0,\dots,\hat g_K\in\mathbb R^d$ и симметрических $r$-линейных отображений $\hat D_{r,k}:(\mathbb R^d)^r\to\mathbb R^d$, $1\le r\le R$, $0\le k\le K$, существует задача~\eqref{eq:fde} с правой частью $f$, аналитической в $(0,\infty)\times\mathbb R^d$, решение которой аналитично при $t>0$ и для которой $g_k=\hat g_k$ и $D_{r,k}=\hat D_{r,k}$ при указанных $r,k$.
\end{lemma}
\begin{proof}
Положим $\eta(t)=\sum_{k\le K}\hat g_k(t-t_n)^k/k!$, $y(t)=y_0+\frac{1}{\G(\alpha)}\int_0^t(t-s)^{\alpha-1}\eta(s)\dd s$; функция $y$ аналитична при $t>0$ и непрерывна на $[0,\infty)$. Определим
\[
f(t,u)=\eta(t)+\sum_{r=1}^{R}\sum_{k=0}^{K}\frac{(t-t_n)^k}{k!\,r!}\,\hat D_{r,k}\bigl(u-y(t),\dots,u-y(t)\bigr).
\]
Тогда $f(t,y(t))=\eta(t)={}^{C}\!D^{\alpha}y(t)$, т.\,е. $y$~--- решение задачи с этой правой частью (единственное в силу локальной липшицевости $f$ по $u$), $g=\eta$, и $\partial_u^rf(t,y(t))=\sum_k\frac{(t-t_n)^k}{k!}\hat D_{r,k}$ при $1\le r\le R$.
\end{proof}

Смысл леммы прост: в начальной точке производные решения определяются правой частью, но во внутренней точке $t_n>0$ значения $y'(t_n),y''(t_n),\dots$ зависят от всей предыстории, и полные производные $g_k$, $D_{r,k}$ можно задать произвольно. Тем самым элементарные дифференциалы~--- независимые полиномы от свободных переменных, и коэффициенты разложения~\eqref{eq:bseries} обязаны обращаться в нуль по отдельности.

\begin{theorem}[необходимость и достаточность условий]\label{thm:necessity}
Пусть $q>0$. (а) Локальная ошибка метода~\eqref{eq:method}, \eqref{eq:rowsum} равна $O(h^q)$ для всех систем с гладкой правой частью тогда и только тогда, когда $\Psi(T)=0$ для всех деревьев с $|T|<q$.
(б) То же для всех скалярных уравнений имеет место тогда и только тогда, когда для каждого одночлена $m$ от атомов
\[
\sum_{T:\ |T|<q,\ \mathfrak m(T)=m}\frac{\Psi(T)}{\sigma(T)}=0,
\]
где $\mathfrak m(T)=\prod_v D_{r_v,k_v}$~--- произведение атомов по стадийным вершинам $T$ (в скалярном случае $F(T)=\mathfrak m(T)$).
\end{theorem}
\begin{proof}
Достаточность следует из теоремы~\ref{thm:bseries} с $P<q$, $P^{+}\ge q$. Докажем необходимость в случае (а). Пусть $T_0$~--- дерево с $|T_0|<q$ и $\Psi(T_0)\ne0$, причём из всех таких деревьев выбрано дерево с наименьшим показателем. Занумеруем стадийные вершины $T_0$ числами $1,\dots,d$ (корень~--- номер $1$) и обозначим через $e_v$ базисные векторы $\mathbb R^d$. Для каждой стадийной вершины $v$ с детьми-стадиями $u_1,\dots,u_r$ и $k$ листьями положим, что $D_{r,k}$ содержит симметризацию отображения $(x_1,\dots,x_r)\mapsto e_v\prod_l\langle e_{u_l},x_l\rangle$ (при $r=0$: $g_k$ содержит слагаемое $e_v$); прочие атомы положим нулевыми. По лемме~\ref{lem:free} существует задача с такими атомами. Первая компонента $F(T)$ отлична от нуля лишь тогда, когда корень $T$ имеет ту же пару $(r,k)$, что и корень $T_0$, а дети-стадии корня $T$ (как мультимножество поддеревьев) переходят в детей-стадий корня $T_0$ с сохранением этого свойства; рекурсивно это означает $T\cong T_0$, и тогда $F(T_0)_1=\prod_v(\text{положительные комбинаторные множители})>0$. Следовательно, первая компонента коэффициента при $h^{|T_0|}$ в~\eqref{eq:bseries} равна $\Psi(T_0)F(T_0)_1/\sigma(T_0)\ne0$ (деревья с тем же показателем, но не изоморфные $T_0$, в первую компоненту не вносят вклада даже при рациональном $\alpha$, когда показатели разных пар $(N,K)$ совпадают), а все члены с меньшими показателями обращаются в нуль по выбору $T_0$. Значит, $\delta_n\ne O(h^q)$ для этой задачи. В случае (б) атомы~--- свободные вещественные переменные (лемма~\ref{lem:free} при $d=1$), одночлены от них линейно независимы как функции, и рассуждение то же.
\end{proof}

\begin{remark}
Теорема~\ref{thm:necessity} восполняет звено, неявно использованное в доказательствах барьеров в~\cite{PrikBarriers}: несовместность системы условий означает невозможность показателя лишь при условии, что условия \emph{необходимы}. Для классических методов Рунге--Кутты необходимость обеспечивается независимостью элементарных дифференциалов Бутчера; здесь её роль играет лемма~\ref{lem:free}.
\end{remark}

\subsection{Число условий и скалярный случай}

Число деревьев с $N$ стадийными вершинами и $K$ листьями равно числу корневых деревьев с $N+K$ вершинами и $K$ листьями (последовательность A055277 в~OEIS). При $\alpha=1$ уровни с одинаковой суммой $N+K$ сливаются, и получаются классические числа условий порядка $1,1,2,4,9,20,48,\dots$ При $\alpha<1$ уровни расходятся по разным степеням $h$, и число условий, которые нужно обнулить для достижения данного показателя, зависит от $\alpha$. Сводка приведена в табл.~\ref{tab:counts}.

\begin{table}[ht]
\centering\small
\caption{Число условий порядка на уровне $h^{K+N\alpha}$: для систем / для скалярных уравнений.}\label{tab:counts}
\begin{tabular}{c|cccccc}
\toprule
$N\backslash K$ & 0 & 1 & 2 & 3 & 4 & 5\\
\midrule
1 & 1/1 & 1/1 & 1/1 & 1/1 & 1/1 & 1/1\\
2 & --- & 1/1 & 2/2 & 3/3 & 4/4 & 5/5\\
3 & --- & 1/1 & 4/3 & 8/6 & 14/10 & 21/15\\
4 & --- & 1/1 & 6/3 & 18/8 & 39/16 & 72/28\\
5 & --- & 1/1 & 9/3 & 35/9 & 97/21 & 214/42\\
6 & --- & 1/1 & 12/3 & 62/9 & 212/23 & 563/52\\
\bottomrule
\end{tabular}
\end{table}


При $K=1$ на каждом уровне имеется единственное дерево~--- цепочка $B_{N-1}$. При $K=2$ число деревьев для систем равно $\lfloor(N+1)^2/4\rfloor$, а для скалярных уравнений при $N\ge3$ равно трём: деревья с одной вставкой $J'$ в цепочку (вставки в разные места дают один одночлен $J^{N-2}J'g_1$), цепочка $b^{\top}A^{N-2}d_2$ при $J^{N-1}g_2$ и деревья с одной вставкой $f''$. Скалярная редукция впервые проявляется на уровне $h^{2+3\alpha}$, где деревья $J(J'g_1)$ и $J'(Jg_1)$ с весами $b^{\top}A(c\cdot d_1)$ и $b^{\top}(c\cdot Ad_1)$ дают одну и ту же величину $JJ'g_1$; это объясняет различие \guillemotleft векторного\guillemotright\ и \guillemotleft скалярного\guillemotright\ вариантов ядра условий, отмеченное в~\cite{PrikBarriers}.

Полный список деревьев с $N+K\le6$ (37 условий для систем) с весами и элементарными дифференциалами, сформированный программой, приведён в репозитории (файл \texttt{paper2/data/conditions\_N+K\_le\_7.md}); условия, нужные для пятистадийных методов, выписаны в табл.~\ref{tab:S} п.~\ref{sec:s5}.
